% Generated draft: requires mathematical and source review.
\documentclass[11pt]{article}
\usepackage[margin=25mm]{geometry}
\usepackage{fontspec}
\setmainfont{Latin Modern Roman}
\usepackage{amsmath,amssymb,mathtools,unicode-math}
\setmathfont{Latin Modern Math}
\usepackage{enumitem,needspace,xurl,hyperref}
\setlength{\parindent}{0pt}
\setlength{\parskip}{5pt}
\setlength{\emergencystretch}{3em}
\newcommand{\N}{\mathbb{N}}
\newcommand{\Z}{\mathbb{Z}}
\newcommand{\Q}{\mathbb{Q}}
\newcommand{\R}{\mathbb{R}}
\newcommand{\C}{\mathbb{C}}
\newcommand{\F}{\mathbb{F}}
\newcommand{\A}{\mathbb{A}}
\newcommand{\PP}{\mathbb{P}}
\newcommand{\E}{\mathbb{E}}
\newcommand{\eps}{\varepsilon}
\newcommand{\dd}{\,\mathrm{d}}
\begin{document}
\section*{UnsolvedMath ID: 30003428}
\textbf{Source problem code:} catalogue entry
\subsection{Definitions and mathematical statement}
For a parameter \(H\in(0,1)\), a normalized fractional Brownian motion \((B_t^H)_{t\geq 0}\) is a centered real Gaussian process with continuous paths and covariance
\[
\mathbb{E}[B_t^H B_s^H]
=\frac12\bigl(t^{2H}+s^{2H}-|t-s|^{2H}\bigr),
\qquad s,t\geq0.
\]
Its increments satisfy
\[
\mathbb{E}\bigl[(B_{t+h}^H-B_t^H)^2\bigr]=h^{2H},
\qquad t\geq0,\quad h>0.
\]
Here \(H\) is the Hurst parameter. The regime \(0<H<\tfrac12\) is conventionally called rough in volatility modeling; the supplied problem says only that \(H\) is small and supplies neither a numerical value nor a precise range.

\textbf{Research question.} What mathematical or market-microstructural mechanisms explain the apparently universal roughness of financial volatility, commonly modeled using fractional Brownian motion with a small Hurst parameter \(H\)? In this question, apparent universality is an empirical motivation, not an established assumption or a mathematically quantified conclusion.

\textbf{Unresolved formulation.} The source does not specify the volatility observable, a class of market models, assumptions on market dynamics, the relevant time scales, a limiting procedure, or a criterion for agreement with fractional Brownian motion. In particular, it does not specify whether the fractional Brownian model describes volatility fluctuations, log-volatility, or another derived quantity. The Gaussian process defined above supplies a precise reference process, but does not itself identify a financial volatility process or a proposed mechanism. The source therefore poses an explanatory research question rather than a fully quantified conjecture. A precise theorem would require these choices and an explicit meaning of universality across markets; no such choices are imposed here.
\subsection{Short English statement}
Explain why financial volatility appears rough across markets and why models involving fractional Brownian motion with a small Hurst parameter fit this behavior. The source leaves the observable, mechanisms, scales, and meaning of universality unspecified, so this remains an explanatory research question requiring further mathematical formulation.
\subsection{Sources}
\begin{itemize}
\item Supplied local UnsolvedMath registry record: permanent ID \(30003428\), title \emph{Market Origins of Rough Volatility}, collection \emph{Oberwolfach Reports Open Problems}, exact source problem code \texttt{catalogue entry}. No individual proposer is supplied; the proposed year is recorded as 2017.
\item Supplied report citation: \emph{Mathematics of Quantitative Finance (2018)}, Oberwolfach Reports, DOI 10.4171/owr/2017/13, pp.~683--769. \url{https://doi.org/10.4171/owr/2017/13}. The supplied background describes Oberwolfach Reports as documenting mathematical workshops and their open-problem discussions. No exact problem page or named speaker is supplied. The citation's 2018 designation and the DOI's 2017 component are retained without resolving their relationship.
\item Supplied literature reference, labeled ``Benzaquen--Bouchaud--Mastromatteo, microstructural foundations'': \url{https://arxiv.org/abs/1609.05177}. The label and link are retained as supplied; their correspondence has not been independently verified.
\item Supplied literature reference, labeled ``Gatheral--Jaisson--Rosenbaum, rough volatility evidence'': \url{https://arxiv.org/abs/1702.02777}. The label and link are retained as supplied; their correspondence has not been independently verified.
\item Literature assessment dated \texttt{2026-08-21}: the supplied curation labels the problem open and states that microstructural and statistical models generate rough volatility and reproduce several stylized facts, while a universal origin of roughness across markets is not established. This is supplied dated curation derived from the OWR record, not an independently verified assessment; time-sensitive status and cited claims require independent verification.
\item The registry supplies statement status \texttt{corrected\_verified}, explained there as preservation of mathematical content while repairing notation, typography, and sentence structure. This is a supplied editorial designation, not independent mathematical or source verification of this draft.
\end{itemize}
\end{document}
